% REV_P0 2026-11-24：由 chapter_02_background_rewrite_v1_20261122.tex 派生；V1保持不变。
% 本轮已核对CellCharter原文正文、Methods、Fig.1、Extended Data Fig.1--2、10页Supplementary Information
% 及官方仓库CSOgroup/cellcharter_analyses（commit fd63ac5）；旧的“未核对补充材料”声明已改写。
% 第三篇CellCharter毕业论文第二章相关研究（2026-11-22）新版首稿
% 不覆盖历史chapter_02_background_integrated.tex；正文依赖主模板提供\cite。
% 参考文献条目在 thesis_cellcharter_verified_references_20261122.bib。
% 核查范围：已有项目11月11日文献证据矩阵 + 官网2026-11-22核对的BANKSY、BayesSpace、Leiden和Wu原文。
% 未逐项核对CellCharter的所有Supplement尺度消融，不能由“未报告”推断“没有前人做过”。
\section{研究背景与相关工作}
\label{sec:cc-background-rewrite-v1}

\subsection{空间转录组与组织空间域识别的研究需求}
\label{subsec:cc-st-background-v1}

基因表达谱为理解组织细胞状态、发育分化及疾病相关微环境提供了分子证据，但在不保留空间位置信息的测序设计中，细胞或组织区域间的原始相邻关系较难直接恢复。空间转录组测序通过将表达读数与组织中的采样坐标关联，使研究者能够同时考察表达相似性和空间结构。空间域识别（spatial domain identification）旨在将具有相近分子或空间特征的观测单位归入不同区域，所获得的聚类结果可用于探索组织分层、邻近细胞相互作用与肿瘤微环境等问题。

空间域聚类与一般表达聚类既有联系，也存在评价上的差别。对于表达聚类，两个来自相隔较远组织位置的spot可以因为表达相似而被分在同一个簇；对于空间域识别，同一簇在空间上是否连续、簇间交界能否对应真实组织结构，同样是具有解释价值的问题。反之，高度强调空间连续性也可能使小范围结构与清晰层间交界在模型输出中变得不易区分。因此，较高的总体聚类一致性并不自动意味着准确的空间边界定位，二者应有相互独立的评价证据。

人类背外侧前额叶皮层（dorsolateral prefrontal cortex，DLPFC）具有可观察的皮层分层结构，是空间域方法常用的真实数据场景。Maynard等发表了包含3位供体、12张Visium切片的DLPFC空间表达资料，并开展组织分层和表达结构研究\cite{maynard2021}。这些数据可用于检验模型是否能够恢复已有参考组织层，但应特别注意其采样结构：来自同一供体的不同切片并非彼此独立的生物学重复。此外，组织层的参考标注可能受到切片质量、标注来源和层间过渡区域的影响，不应把所有参考标签视为完美且无误差的绝对真值。

本研究除DLPFC外，还在另一种组织来源中评价方法流程。Wu等发布了人类乳腺癌的单细胞和空间分辨图谱，其公开数据中包含六例用于Visium空间分析的乳腺癌样本及相关空间元数据\cite{wu2021}。不同于皮层的层级结构，乳腺癌组织区域具有另一套病理标注语义，因此跨数据集比较的重点应是方法在新组织队列上的表现和患者异质性，而不是将DLPFC的ARI绝对值与乳腺癌的ARI绝对值直接比较后推断某一种组织更易分析。

\subsection{表达降维与空间信息的引入}
\label{subsec:cc-expression-v1}

空间表达矩阵通常维数较高，不同基因在测序深度、零计数比例和技术噪声方面存在差异。直接在全部基因上进行聚类，可能使计算成本和不稳定性同时增加。scVI（single-cell variational inference）是既有的概率深度生成建模框架，可从表达计数中学习适合下游分析的低维潜变量\cite{lopez2018}。本研究使用scVI作为已发表的表达表征工具，并将其5维潜变量作为单spot的基础特征，不涉及设计新的变分推断目标或神经网络结构。

只使用单spot表达，可能无法利用临近区域的组织上下文。空间聚类因而常将表达表示与组织坐标结合：一类方法首先通过表达构建低维表示，再从空间邻域聚合相邻点特征；另一类方法直接利用空间图或原生空间统计模型完成联合聚类。在具体实现上，不同方法可能使用空间最近邻、固定物理网格、像素坐标、阵列坐标或空间先验。它们所使用的空间结构虽然来源于同一组织，却并非一定是可互换的模型输入。因此，比较不同算法时应在统一样本和评价mask的前提下保留方法自身的输入定义，并避免宣称所有算法在完全相同特征与计算预算下受到控制。

SEDR（spatially embedded deep representation）等方法通过深度表示学习把空间联系纳入聚类过程。Xu等提出的SEDR框架包含掩码自监督表示学习与空间图表示等环节，说明“把空间信息引入表达表示”并不是单一模型的特有设计\cite{xu2024}。对本研究而言，其后续DLPFC参考标注的部分整理文件与SEDR相关公开分析资料有数据层面的联系，但标注文件存在于SEDR资料中不等于这些标签由SEDR模型生成，更不意味着标签由本研究重新人工标记。使用参考标签进行方法评价，应另行追溯其原始来源和加工过程。

\subsection{CellCharter及多尺度邻域表征}
\label{subsec:cc-cellcharter-v1}

CellCharter是由Varrone等提出的成熟空间生态位识别、表征及比较框架\cite{varrone2024}。其公开软件支持以已有表示为输入，对不同层数的空间邻域进行特征聚合，并将中心观测单位及多层邻域特征共同用于下游分析。官方软件接口提供聚合层数等参数，使研究者可以基于同一初始表达嵌入构建多尺度表示。本研究采用其公开实现及相应聚类功能，不将该邻域聚合算子本身表述为新提出的方法。

对于某一切片的$N$个spot，假设每个spot的scVI基础表示为5维，则可以把相应的CellCharter输出按顺序写为
\begin{equation}
H=[L_0\mid L_1\mid L_2\mid L_3\mid L_4]
\in\mathbb{R}^{N\times25},\quad L_\ell\in\mathbb{R}^{N\times5}.
\end{equation}
其中$L_0$对应原始基础表示，其余块与软件输出的不同邻域层相对应。在本项目中，该分块结构允许从同一25维冻结缓存中提取5D、10D、15D、20D、25D累计组合，也允许构造不同但维度相同的15D组合。后者适合检查“使用哪一层邻域块”是否与整体聚类和局部组织边界表现相关，但不能自动分离输入尺度、维数、空间拓扑和聚类数变化的全部效应。

多尺度空间信息的构建并非只有CellCharter一种途径。例如，m2ST的研究已讨论多尺度图表示及尺度信息的融合问题\cite{zhang2025}。已有工作还将不同尺度的贡献作为可调对象，说明仅提出“更深邻域可能有用”或者“按不同尺度拼接特征”，不足以建立新的核心算法贡献。本研究因此采用实证评价的定位：明确固定不同维度、等维替换、K选择规则和评价指标，关注在真实数据中哪些结论相对稳定、哪些条件会出现方向相反的结果。

\subsubsection{原作者在DLPFC上已完成的工作}
\label{subsubsec:cc-prior-dlpfc-rev-p0}
CellCharter原文已使用本研究所用的同一套12张DLPFC Visium切片进行基准比较\cite{varrone2024}，与本研究直接相关的内容如下：
\begin{enumerate}
  \item \textbf{调参与测试划分。}原文Methods中“Benchmarking comparison of CellCharter with existing tools”一节规定：每位供体随机选1张切片（151507、151672、151673）组成调参集，用网格搜索、每组参数5次重复、取调参集平均ARI最大者确定超参数；其余9张切片作为测试集，每张重复10次，仅以ARI评价，未标注spot参与聚类、不参与评分。Extended Data Fig.~1e、1f对比了“在同一数据上调参和测试”与“调参集与测试集分离”两种做法。
  \item \textbf{邻域层数已被调参。}Supplementary Information“Hyperparameter tuning: grid search”一节列出CellCharter的候选值：高变基因$\{500,1000,2000,5000\}$，选5000；潜变量维度$\{5,10,15\}$，选5；邻域层数$\{2,3,4\}$，选4。官方仓库\texttt{src/benchmarking/CellCharter/tuning.py}与\texttt{ARI\_CellCharter\_tuning.csv}给出这一网格的逐格ARI\cite{cellcharteranalyses2024}。因此在5000高变基因、5维潜变量条件下，2--4层邻域对ARI的影响已有原作者在3张调参切片上的评价。
  \item \textbf{单切片与联合聚类。}原作者统一要求各方法输出7个簇（Results正文）；但官方单切片脚本\texttt{individual.py}对151669--151672设$K=5$、其余设$K=7$，与本研究的Oracle-K相同。原文的主要性能结论来自9张测试切片的\textbf{联合聚类}（Fig.~1d、1e；Extended Data Fig.~1h）。在\textbf{单切片}聚类中，原文报告STAGATE表现最好（Extended Data Fig.~1g）\cite{dong2022stagate}。
  \item \textbf{无标签选K。}CellCharter的选K规则是：对每个候选K重复拟合GMM，按K、$K-1$、$K+1$的Fowlkes--Mallows指数\cite{fowlkes1983}判断稳定性（原文Fig.~1a及Methods“Cluster stability”）。在DLPFC上，原作者对全部12张切片\textbf{联合}运行该规则（候选2--20，10次重复），稳定性峰值为$n=9$。同一设置下BIC的全局最小值在$n=19$，AIC与负对数似然没有清晰拐点（Fig.~1h；Extended Data Fig.~2b；Supplementary Information中关于cluster stability与AIC/BIC/NLL比较的一节）。
\end{enumerate}

据此，本研究与原作者工作的重叠和差异如下。重叠部分包括：25D配置、Oracle-K设置、ARI评价，以及2--4层邻域对ARI影响的比较方向。本研究另外完成的部分包括：0层与1层累计特征和非连续等维15D组合；在全部12张切片上每切片10个GMM种子的重复结构；预设固定K网格（E3c）；\textbf{逐切片}的稳定性AutoK及其事后ARI/NMI；严格与容差空间边界指标；BANKSY与BayesSpace双K臂对照；乳腺癌外部队列。原作者的无标签选K是在联合样本上作示例，而不是逐切片的标签评分。本研究不作“首次发现聚合尺度影响聚类性能”的声明；对于“既有研究未评估空间边界”，本研究只核对了CellCharter原文与官方仓库，不能推广到全部文献。

\subsection{成熟聚类基线与不同K选择方式}
\label{subsec:cc-baselines-background-v1}

非空间GMM是常见的概率混合聚类思路，能够将低维表达表示划分成若干统计成分；但在本项目的非空间基线中，其输入不包含显式组织坐标。以BIC等似然准则为GMM选K时，结果容易偏向更多簇；CellCharter原作者在DLPFC联合数据上也观察到BIC建议的簇数（$n=19$）明显多于稳定性峰值（$n=9$）\cite{varrone2024}。Leiden是一种图社区发现算法，Traag等针对图社区划分提出了保证社区连通性的改进方法\cite{traag2019}。当Leiden用于表达近邻图时，其邻居关系主要由表达表征建立，并不自动等于组织中的物理空间六邻接。二者作为非空间对照，有助于描述额外引入空间表示后可能出现的评价差异。

BANKSY通过融合自身表达及邻域信息构建空间增强的表示，并可结合下游聚类完成空间域识别；其原始研究还展示了组织域分割和细胞类型分析等任务\cite{singhal2024}。BayesSpace则采用低维表达与空间先验相结合的统计模型进行空间聚类，并可支持更精细的空间区域推断\cite{zhao2021}。这两种方法依赖不同的原生空间表示和参数选择程序，不能因为都引入了空间信息就认定它们的计算量或聚类数选择机制等价。

除了模型本身，预测簇数$K$的确定方式也会显著影响结果解释。若直接以人工参考标签的类别数量设置$K$，这实际上使用了标签类别数先验，本文称之为\textit{Oracle-K}。即使模型拟合过程没有读入每个spot的参考层标签，该类别数量仍会影响最终输出。相对地，\textit{K-blind}表示在决定K或无监督分辨率时不使用人工参考层类别信息，例如依据稳定性或BIC的预设规则，或采用固定分辨率。两者应分表呈现，不能把Oracle-K的表现误写成完全无标签的真实部署能力。

不同原生方法也存在选K判据的适用条件。以BayesSpace为例，其\texttt{qTune}给出不同候选q的平均伪对数似然，官方使用说明建议\textbf{在\texttt{qPlot}曲线的拐点（elbow）附近选q}\cite{bayesspacevignette}。本研究冻结的K-blind规则是取平均伪对数似然最大的q，\textbf{与官方建议不一致}；该规则在本研究数据上频繁选中搜索上界（见第\ref{subsec:cellcharter-e5a-v1}节）。原冻结结果全部保留，不在观察结果后更换判据。这一问题属于对照协议的系统性缺陷，会降低BayesSpace在K-blind条件下的可比性，不能当作随机波动处理。

\subsection{聚类指标、空间边界与参考标注的评价问题}
\label{subsec:cc-evaluation-background-v1}

ARI（adjusted Rand index）和NMI（normalized mutual information）是比较预测聚类与参考标签的一类常用指标，二者均不依赖预测簇名称本身的编号。ARI具有机会校正的随机划分参照，NMI衡量两种分区之间的共享信息，但两者本质上比较的是观测单位所属的类别划分，而不是直接测量组织空间边界定位。如果一个算法把相邻真实层间交界略微移开，即使总体ARI没有明显变化，该偏移仍可能改变组织学边界检出的结果；反过来，较少的预测边界也可能同时减少误报和漏检，综合F1需要结合基础计数理解。

基于此，本研究在Visium阵列六邻接图上设置严格edge-F1，以相邻两个spot的参考层是否不同作为真实边界、预测簇是否不同作为预测边界；同时记录TP、FP、FN以区分正确检出、误报和漏检。考虑到严格网格边匹配对少量位置偏移敏感，另将1-hop和2-hop容差作为明确标注的\textbf{事后敏感性分析}，在有效边的线图中使用一对一匹配，避免同一参考边界被重复计算。宽松容差可以回答“邻近交界是否被预测”，不能代替严格定位结果，更不能根据哪个容差指标更有利再调整主评价目标。

已有针对空间聚类评价的研究也提示，在数据处理、标签质量和模型调节条件不同的情况下，性能排序具有不确定性。例如，SACCELERATOR围绕空间感知聚类的评价与专家参与参考机制提出了相关讨论\cite{sun2026}。本研究不将全部参考标签等同于绝对无误差的生物学真值；对于DLPFC皮层层级和乳腺癌病理区域，必须分别说明其标签粒度、数据来源、有效mask和不可直接互换的语义。

对于DLPFC，12张切片来源于3位供体，同一供体内的空间邻近切片不是12份独立受试者；对于外部队列，则有6名不同患者。技术随机种子只能描述算法初始化或拟合波动，不能增加生物学样本量。因此，本研究优先报告按切片/患者分层的效应量、方向一致性、负结果及解释限制，不以spot、边或seed数目制造显著性。该处理也是将整体ARI/NMI与局部边界评价放在同一证据框架中的必要前提。

\subsection{研究问题与本文工作定位}
\label{subsec:cc-questions-background-v1}

结合以上方法特征与评价限制，本研究使用成熟CellCharter流程形成可复核的系统实证评价，具体围绕如下研究问题：
\begin{enumerate}
  \item 在固定数据、GMM与Oracle-K条件下，累计邻域深度和等维邻域块替换如何影响DLPFC整体聚类指标？这一影响对技术随机种子和切片/供体是否稳定？
  \item 当K通过固定网格或不读取人工标签的AutoK规则决定时，尺度比较结果是否发生变化？选K边界和不同实际簇数应如何呈现？
  \item 严格空间边界F1与1-hop、2-hop邻近容差F1是否给出相同尺度排序？误报与漏检是否存在解释上的权衡？
  \item 在同一DLPFC数据上与非空间及成熟空间方法比较时，K-blind/Oracle-K和原生预处理差异会如何限制结论？
  \item 将同类流程重新用于六名乳腺癌患者时，邻域尺度的描述性收益是否具有患者间的一致性；哪些失败与负结果不能被宏平均遮盖？
\end{enumerate}

本文不把这些问题包装成新基础算法的提出和性能保证。研究工作包括真实DLPFC数据的多尺度与选K实验、冻结预测的空间边界和容差评价、成熟基线的双K策略对照、独立外部患者队列的重复分析，以及对所有原始预测、来源SHA和评价脚本的归档。现有实验证据支持的是在所选样本与协议下对方法行为的描述和评价；是否对其他组织、平台或更大患者群体成立，仍需要额外证据，而不能由本章直接推断。
